% frontmatter/abstract-en.tex — English abstract
\chapter*{Abstract}
\addcontentsline{toc}{chapter}{Abstract (English)}

\noindent\textbf{Title:}\ \thesistitleEN\\
\textbf{By:}\ \authoroneEN, \authortwoEN, \authorthreeEN\\
\textbf{Advisor:}\ \advisorEN, Ph.D.\\
\textbf{Academic Year:}\ \academicyearEN

\vspace{0.6cm}

YouTube hosts more than 2.5 billion monthly active users and is the dominant
distribution channel for Thai content creators. Predicting whether a newly
posted video will go \emph{viral} is therefore commercially and academically
valuable. This study formulates the task as binary classification of Thai
videos using a \emph{per-channel virality index} ($\mathit{VI}$) that
aggregates the engagement rate, like ratio, and comment ratio after a
log-transform and a within-channel z-score, with the positive class taken
as the top decile per channel. The per-channel construction explicitly
removes channel size as a confounder.

The dataset comprises 23{,}431 videos from 82 channels collected through the
YouTube Data API v3, with a 10.16\% positive class rate. The split policy is
\emph{channel-grouped and time-aware}: the test set is the latest 15\% of
videos by \texttt{published\_at}, with no channel shared between train and
validation. Class imbalance is handled \emph{only on the training fold}
(3:1 undersampling), so validation and test retain the natural distribution.

Three Thai transformer encoders of differing scale and architecture are
compared end-to-end: \textbf{WangchanBERTa} (105M, RoBERTa SentencePiece),
\textbf{PhayaThaiBERT} (278M, RoBERTa), and \textbf{XLM-RoBERTa-large}
(560M, multilingual). All three are full-fine-tuned on cloud GPUs (NVIDIA
T4/A10G via Hugging Face Jobs). Classical baselines (Logistic Regression,
LightGBM, XGBoost) are trained over 27 structured features fused with
character-N-gram TF-IDF and Wisesight-Sentiment scores. A hybrid model
(MLP + LightGBM head) operates on title embeddings concatenated with
structured features. Finally, a \emph{Platt-calibrated stacking ensemble}
fits a logistic-regression meta-learner over the validation-fold
probabilities of fifteen base models.

On the held-out time-aware test set, the calibrated stacking ensemble
attains a test \textbf{ROC-AUC of 0.6914} with a 1{,}000-sample bootstrap
95\% confidence interval of \textbf{[0.6625, 0.7174]} and an Expected
Calibration Error of 0.016. This is significantly higher than the strongest
single model, LightGBM-Structured+TF-IDF (0.6728; McNemar
$p = 6.9\!\times\!10^{-53}$). Among single transformer encoders, PhayaThaiBERT
performs best (0.6451), followed by WangchanBERTa (0.6402) and
XLM-RoBERTa-large (0.5450). Cochran's Q across the three encoders yields
$Q = 612.69$, $p = 9.0\!\times\!10^{-134}$ ($df = 2$), confirming that the
encoders' error patterns differ significantly. Per-bucket evaluation shows
the ensemble generalises better on large channels ($\geq$1M subscribers,
ROC-AUC = 0.7459) than on mid-size channels (100k--1M, 0.6273).

We report an honest \emph{negative finding}: feeding PhayaThaiBERT a
multi-field text (\texttt{title + description + channel}) \emph{degraded}
ROC-AUC from 0.6451 to 0.5455, reflecting the boilerplate nature of Thai
YouTube descriptions. Explainability analyses (SHAP, LIME, attention
rollout) show that the dominant predictive signal lives in structured
features (title length, emoji density, channel subscriber count) and that
language signals contribute to the ensemble through diversity rather than
through any single best encoder.

\noindent\textbf{Keywords:} YouTube virality prediction, Thai transformer
language models, WangchanBERTa, PhayaThaiBERT, XLM-RoBERTa, stacking
ensemble, McNemar's test, Cochran's Q, probability calibration, SHAP, LIME.

\newpage
